\bta{2021}


\begin{enumerate}
	%\renewcommand{\labelenumi}{\arabic{enumi}.}
	% A(\Alph) a(\alph) I(\Roman) i(\roman) 1(\arabic)
	%设定全局标号series=example	%引用全局变量resume=example
	%[topsep=-0.3em,parsep=-0.3em,itemsep=-0.3em,partopsep=-0.3em]
	%可使用leftmargin调整列表环境左边的空白长度 [leftmargin=0em]
	\item
	（30分） 能量为$  E $的粒子, 从负无穷远处沿着 $x$ 轴入射, 势场为
	\[ 
	V(x)=\left\{\begin{aligned}&0, &  x<0 \\ &V_{0}, & 0<x<a \\ &\infty, &  \quad x>a\end{aligned}\right.
	 \]
	其中 $E>V_{0}>0$, 求全空间的波函数 $\psi(x)$ 。


\banswer{
	令
	\[
		k=\frac{\sqrt{2mE}}{\hbar},\qquad
		q=\frac{\sqrt{2m(E-V_0)}}{\hbar}.
	\]
	在 $x>a$ 处有无限高势垒，故波函数为零；在 $0<x<a$ 内选择自动满足
	$\psi(a)=0$ 的形式。取入射振幅为 $1$，全空间波函数可写为
	\[
		\boxed{\psi(x)=
		\begin{cases}
		e^{ikx}+r e^{-ikx},&x<0,\\[3pt]
		C\sin[q(a-x)],&0<x<a,\\[3pt]
		0,&x>a,
		\end{cases}}
	\]
	其中在 $x=0$ 的连续匹配条件为
	\[
	 1+r=C\sin(qa),\qquad ik(1-r)=-Cq\cos(qa).
	\]
	由第二式加上 $ik$ 倍第一式，得到
	$2ik=C[ik\sin(qa)-q\cos(qa)]$，从而
	\[
		\boxed{r=\frac{ik+q\cot(qa)}{ik-q\cot(qa)}},\qquad
		\boxed{C=\frac{2ik}{ik\sin(qa)-q\cos(qa)}}.
	\]
	对实数 $k,q$ 有 $|r|=1$，符合右侧硬墙使总反射率为 $1$。
}

	
\item 
（30分）已知 $H$ 的两个本征态 $\left|u_{1}\right\rangle,\left|u_{2}\right\rangle$ 分别对应能量 $E_{1}, E_{2}$。力学量 $A$ 的本征态为
$\left|\phi_{1}\right\rangle=\frac{\left|u_{1}\right\rangle+\left|u_{2}\right\rangle}{\sqrt{2}}$ 和 $\left|\phi_{2}\right\rangle=\frac{\left|u_{1}\right\rangle-\left|u_{2}\right\rangle}{\sqrt{2}}$，分别对应本征值 $a_{1}, a_{2}$。$t=0$ 时粒子处于 $\left|\phi_{1}\right\rangle$ 态。求：
\begin{enumerate}
	%\renewcommand{\labelenumi}{\arabic{enumi}.}
	% A(\Alph) a(\alph) I(\Roman) i(\roman) 1(\arabic)
	%设定全局标号series=example	%引用全局变量resume=example
	%[topsep=-0.3em,parsep=-0.3em,itemsep=-0.3em,partopsep=-0.3em]
	%可使用leftmargin调整列表环境左边的空白长度 [leftmargin=0em]
	\item
	$t$ 时刻的态 $|\psi(t)\rangle$;
	\item
	 $t$ 时刻力学量 $A$ 的期望值 $\langle A\rangle(t)$。
	
	
	
\end{enumerate}

\banswer{
	\begin{enumerate}
		\item
		\[
			\boxed{|\psi(t)\rangle=\frac1{\sqrt2}
			\left(e^{-iE_1t/\hbar}|u_1\rangle
			+e^{-iE_2t/\hbar}|u_2\rangle\right)}.
		\]
		把它改写到 $A$ 的本征基中：
		\[
			|\psi(t)\rangle=\frac12\left[
			(e^{-iE_1t/\hbar}+e^{-iE_2t/\hbar})|\phi_1\rangle
			+(e^{-iE_1t/\hbar}-e^{-iE_2t/\hbar})|\phi_2\rangle\right].
		\]

		\item 令 $\Delta E=E_2-E_1$，则测得 $a_1,a_2$ 的概率分别为
		\[
			P_1(t)=\cos^2\frac{\Delta Et}{2\hbar},\qquad
			P_2(t)=\sin^2\frac{\Delta Et}{2\hbar}.
		\]
		所以
		\[
			\boxed{\langle A\rangle(t)
			=a_1\cos^2\frac{\Delta Et}{2\hbar}
			+a_2\sin^2\frac{\Delta Et}{2\hbar}}
		\]
		或等价地
		\[
			\langle A\rangle(t)=\frac{a_1+a_2}{2}
			+\frac{a_1-a_2}{2}\cos\frac{\Delta Et}{\hbar}.
		\]
	\end{enumerate}
}

\item 
（30分）
对于自旋 $\frac{1}{2}$ 的体系, 已知哈密顿量 $H=S_{x}+S_{y}$。
\begin{enumerate}
	%\renewcommand{\labelenumi}{\arabic{enumi}.}
	% A(\Alph) a(\alph) I(\Roman) i(\roman) 1(\arabic)
	%设定全局标号series=example	%引用全局变量resume=example
	%[topsep=-0.3em,parsep=-0.3em,itemsep=-0.3em,partopsep=-0.3em]
	%可使用leftmargin调整列表环境左边的空白长度 [leftmargin=0em]
	\item
	求 $H$ 的本征值和本征态;
	\item
	求在 (1) 问的低能级态 $\left|\psi_{-}\right\rangle$下测到 $S_{z}=+\frac{\hbar}{2}$ 的概率;
	\item 
	设 $t=0$ 时体系初态为 $\left|\psi_{-}\right\rangle $。$ t>0$ 时加入 $z$ 方向磁场 $\vec{B}=(0,0, B)$, 求 $t$ 时刻粒子处于初态的概 率。
	
\end{enumerate}

\banswer{
	\begin{enumerate}
		\item 在 $S_z$ 表象中
		\[
			H=\frac{\hbar}{2}\begin{pmatrix}0&1-i\\1+i&0\end{pmatrix}.
		\]
		特征方程为 $E^2-\hbar^2/2=0$。
		若本征矢为 $(u,v)^T$，首行给出
		$v/u=2E/[\hbar(1-i)]=\pm e^{i\pi/4}$；两分量模长相等，归一化后得
		\[
			\boxed{E_\pm=\pm\frac{\hbar}{\sqrt2}},\qquad
			\boxed{|\psi_\pm\rangle=\frac1{\sqrt2}
			\begin{pmatrix}1\\ \pm e^{i\pi/4}\end{pmatrix}}.
		\]

		\item 低能态 $|\psi_-\rangle$ 的上分量模方为 $1/2$，故
		\[
			\boxed{P(S_z=+\hbar/2)=\frac12}.
		\]

		\item 若磁场的 Zeeman 项为 $-\gamma B S_z$，且原哈密顿量在
		$t>0$ 后仍然存在，则
		\[
			H_{\rm tot}=S_x+S_y-\gamma B S_z.
		\]
		记 $\vec d=(1,1,-\gamma B)$、$D=|\vec d|$，则
		$H_{\rm tot}=(\hbar/2)\vec d\cdot\vec\sigma$，且
		\[
		 U=I\cos\frac{Dt}{2}-i\frac{\vec d\cdot\vec\sigma}{D}\sin\frac{Dt}{2}.
		\]
		初态的 Bloch 矢量为 $\langle\vec\sigma\rangle_-=-(1,1,0)/\sqrt2$，故
		\[
		 \langle\psi_-|U|\psi_-\rangle=\cos\frac{Dt}{2}
		       +i\frac{\sqrt2}{D}\sin\frac{Dt}{2}.
		\]
		取模平方得到
		\[
			\boxed{P_-(t)=1-\frac{(\gamma B)^2}{2+(\gamma B)^2}
			\sin^2\left(\frac{Dt}{2}\right)}.
		\]
		这里系数 $1$ 与 $\gamma B$ 必须采用同一频率单位。若题意是
		$t>0$ 时关闭原来的 $S_x+S_y$、只保留 Zeeman 项，则相应结果为
		\[
			\boxed{P_-(t)=\cos^2\left(\frac{\gamma Bt}{2}\right)}.
		\]
		题干没有说明原哈密顿量是否撤去，以上分别列出两种物理情形。
	\end{enumerate}
}



\newpage
\item 
（30分）
在宽度为$ 2a $的一维无限深势井中存在两个粒子,忽略相互作用。
\begin{enumerate}
	%\renewcommand{\labelenumi}{\arabic{enumi}.}
	% A(\Alph) a(\alph) I(\Roman) i(\roman) 1(\arabic)
	%设定全局标号series=example	%引用全局变量resume=example
	%[topsep=-0.3em,parsep=-0.3em,itemsep=-0.3em,partopsep=-0.3em]
	%可使用leftmargin调整列表环境左边的空白长度 [leftmargin=0em]
	\item
	求能级 $E_{n_{1}, n_{2}}$;
	\item
	若两粒子为自旋 $\frac{1}{2}$ 的全同粒子，求基态与第一激发态的能级简并度以及相应的波函数;
	\item
	若两粒子为自旋 $0$ 的全同玻色子，求基态与第一激发态的能级简并度以及相应的波函数。
	
	
	
\end{enumerate}
\banswer{
	设势阱为 $-a<x<a$，单粒子归一化本征函数记为 $\phi_n(x)$，其能量为
	\[
		\varepsilon_n=\frac{n^2\pi^2\hbar^2}{8ma^2},\qquad n=1,2,\ldots.
	\]
	\begin{enumerate}
		\item 两粒子无相互作用时
		\[
			\boxed{E_{n_1n_2}=\varepsilon_{n_1}+\varepsilon_{n_2}
			=\frac{\pi^2\hbar^2}{8ma^2}(n_1^2+n_2^2)}.
		\]

		\item 对两个全同自旋 $1/2$ 费米子，基态能量为
		\[
			E_{\rm g}=2\varepsilon_1,
		\]
		唯一允许态是
		\[
			\boxed{\phi_1(x_1)\phi_1(x_2)\chi_{00}},
		\]
		其中 $\chi_{00}$ 为自旋单态，故简并度为 $1$。

		第一激发能量为
		\[
			E_{\rm e}=\varepsilon_1+\varepsilon_2=5\varepsilon_1.
		\]
		其全部四个态为
		\[
			\boxed{\frac{\phi_1(x_1)\phi_2(x_2)+\phi_2(x_1)\phi_1(x_2)}
			{\sqrt2}\chi_{00}},
		\]
		以及
		\[
			\boxed{\frac{\phi_1(x_1)\phi_2(x_2)-\phi_2(x_1)\phi_1(x_2)}
			{\sqrt2}\chi_{1m}},\qquad m=1,0,-1.
		\]
		所以第一激发态简并度为 $4$。

		\item 对两个自旋为 $0$ 的全同玻色子，没有额外的自旋简并，空间波函数必须对称。
基态为两个粒子均占据单粒子基态：
\[
\boxed{\Psi_{\rm g}^{\rm B}(x_1,x_2)=\phi_1(x_1)\phi_1(x_2)},
\qquad E_{\rm g}=2\varepsilon_1,\qquad \boxed{g_{\rm g}=1}.
\]
第一激发态为一个粒子占据 $n=1$、另一个占据 $n=2$ 的对称组合：
\[
\boxed{\Psi_{\rm e}^{\rm B}(x_1,x_2)=
\frac{\phi_1(x_1)\phi_2(x_2)+\phi_2(x_1)\phi_1(x_2)}{\sqrt2}}.
\]
由单粒子态的正交归一性，上式已归一化；交换粒子标签不产生新的独立态。
因此
\[
E_{\rm e}=\varepsilon_1+\varepsilon_2=5\varepsilon_1,
\qquad \boxed{g_{\rm e}=1}.
\]
\end{enumerate}
}



\item 
（30分）
氢原子处于基态,受到脉冲电场
$$
\mathscr{E}(t)=\mathscr{E}_{0} \delta(t)
$$
作用, $\mathscr{E}_{0}$ 为常数。已知波尔半径为 $a_{0}$, 矩阵元 $z_{n 1}=\iiint \psi_{n}^{*} z \psi_{1} \mathrm{dV}$。
\begin{enumerate}
	%\renewcommand{\labelenumi}{\arabic{enumi}.}
	% A(\Alph) a(\alph) I(\Roman) i(\roman) 1(\arabic)
	%设定全局标号series=example	%引用全局变量resume=example
	%[topsep=-0.3em,parsep=-0.3em,itemsep=-0.3em,partopsep=-0.3em]
	%可使用leftmargin调整列表环境左边的空白长度 [leftmargin=0em]
	\item
利用选择定则判断粒子会跃迁到哪些激发态上;
	
	
	\item
	计算跃迁到各个激发态的概率;
	
	
	\item
	计算$ t $时刻粒子处于基态的概率。
	
	
	
	
	
	
\end{enumerate}


\banswer{
	取脉冲沿 $z$ 方向。电子与电场的相互作用在符号上可写成
	$H_I(t)=e\mathscr E_0z\,\delta(t)$；改变整体符号不会影响以下概率。
	\begin{enumerate}
		\item $z=r\sqrt{4\pi/3}\,Y_1^0$ 是秩一、$q=0$ 的球张量。从氢原子基态
		$|100\rangle$ 出发，电偶极选择定则为
		\[
			\Delta l=\pm1,\qquad \Delta m=0.
		\]
		因此离散激发态只能是
		\[
			\boxed{|n,1,0\rangle,\qquad n=2,3,\ldots}
		\]
		（完整谱中还包括 $l=1,m=0$ 的连续态）。

		\item 一级含时微扰论给出脉冲后的振幅
		\[
			c_{n10}^{(1)}=-\frac{i e\mathscr E_0}{\hbar}z_{n1}.
		\]
		故跃迁到各个允许离散激发态的概率为
		\[
			\boxed{P_{n10}=\frac{e^2\mathscr E_0^2}{\hbar^2}|z_{n1}|^2,
			\qquad n\ge2}.
		\]

		\item 脉冲结束后各展开系数只积累相位，所以基态概率不再随 $t$ 改变。保留到
		$\mathscr E_0^2$，
		\[
			\boxed{P_{100}(t)=1-\frac{e^2\mathscr E_0^2}{\hbar^2}
			\sum_{\rm excited}|z_{n1}|^2}.
		\]
		若求和包含全部离散与连续激发态，由完备性及
		$\langle100|z|100\rangle=0$，
		\[
			\sum_{\rm excited}|z_{n1}|^2
			=\langle100|z^2|100\rangle=a_0^2,
		\]
		从而
		\[
			\boxed{P_{100}(t)=1-\left(\frac{e\mathscr E_0a_0}{\hbar}\right)^2
			+O(\mathscr E_0^4)}.
		\]
	\end{enumerate}
}


	
\end{enumerate}
